% =====================================================================
% Spin III: The ZBW Mode Spectrum and the Lattice Voronoi-Cell
% Eigenvalue Problem
% Third paper of the CPP spin sub-arc (Spin I, Spin II, Spin III).
%
% CHANGELOG
%   V0   19 Aug 2026  Patch 3234  Initial draft (Session 149), written
%        from the arc's own materials at founder request: the
%        development-notes charter ("Spin III: derive the sub-harmonic
%        spectrum from the 600-cell Voronoi eigenvalues"), the committed
%        eigenvalue script (16 Mar 2026), its data and figure. No Spin
%        III paper previously existed (full-history finding, Patch 3232
%        record). Honest-content ruleset: the rigorous pieces (mode
%        arithmetic, FD verification, scale-bridge identity,
%        discreteness bound) are stated as results; the lattice-cell
%        capstone is reported as OPEN with the committed instrument's
%        boundary-condition mismatch diagnosed and a corrected
%        instrument frozen (OPEN-QM-8). Verify:
%        scripts/3234_verify_spin3_v0.py (10/10).
% =====================================================================
\documentclass[12pt]{article}

\usepackage{amsmath,amssymb,amsthm}
\usepackage{geometry}
\usepackage{hyperref}
\usepackage{booktabs}
\usepackage{enumitem}
\usepackage{parskip}
\geometry{margin=1in}

% ---- corpus macros (matching the Spin I/II conventions) ----
\newcommand{\lP}{l_P}
\newcommand{\SSV}{\text{SSV}}
\newcommand{\rth}{r_{\rm th}}
\newcommand{\rin}{r_{\rm in}}
\newcommand{\rout}{r_{\rm out}}
\newcommand{\rC}{r_C}
\newcommand{\aBohr}{a_0}

\newtheorem{theorem}{Theorem}
\newtheorem{proposition}{Proposition}
\newtheorem{lemma}{Lemma}
\newtheorem{remark}{Remark}

\title{Spin III: The ZBW Mode Spectrum and the Lattice
Voronoi-Cell Eigenvalue Problem\\
\normalsize Version 0 (draft) --- Conscious Point Physics spin sub-arc}
\author{Thomas Lee Abshier, ND\\ Hyperphysics Institute}
\date{19 August 2026}

\begin{document}
\maketitle

\begin{abstract}
This paper is the capstone of the CPP spin sub-arc.  Spin~I established
that the captured Dipole Particle orbital geometry yields the electron
spin angular momentum $L = \hbar/2$ exactly, \emph{given} the radius
condition $\rout = 2\rin$.  Spin~II derived that condition from the ZBW
polarization cloud's standing-wave structure: an open--closed resonator
(node at $r=0$ by CP Exclusion, antinode at the thermal boundary
$\rth = \hbar/2m_ec$) whose Mode~2 has an interior antinode at $\rth/3$
and an interior node at $2\rth/3$ --- ratio $2$, exactly, independent of
every physical constant.  Spin~III's charter, set when Spin~II was
written, is to ground that mode structure in the substrate itself: to
derive the spectrum from the eigenvalue problem of the lattice
Voronoi cell.  This V0 establishes what can currently be established
and reports plainly what cannot.  \textbf{Established here:} the exact
resonator spectrum $k_n = (2n-1)\pi/2\rth$ with finite-difference
verification to $0.01\%$ across eight modes by two independent
constructions; the Mode-2 geometry to $0.1\%$; the scale-bridge
identity $\rin^{\rm (I)}/\rin^{\rm (II)} = 6/[4\alpha(1+\sqrt{2})^2]
= 35.2675$, shown to be an \emph{exact algebraic identity} rather than
a numerical coincidence; and the discreteness bound
$(\lP/\rth)^2 \approx 7\times10^{-45}$, which fixes the division of
labor --- the lattice must supply mode \emph{selection}, not mode
\emph{corrections}.  \textbf{Diagnosed here:} the arc's committed
lattice-cell computation (a graph Laplacian on the 24-cell interior)
solves the closed (Neumann-type) problem rather than the open--closed
one --- the CP-Exclusion node at $r=0$ is never encoded --- so its low
spectrum consists of the four 4D dipole modes and its own Mode-2
diagnostic correctly reports no identification.  \textbf{Registered
here:} the corrected instrument (Dirichlet core, free outer boundary,
isotropy-scored radial profiling), with readings frozen before its
data exists, and the substrate-domain question (24-cell surrogate
versus the true GP-lattice Voronoi cell) put to the founder ---
together \texttt{OPEN-QM-8}.
\end{abstract}

\tableofcontents

% =====================================================================
\section{The arc and this paper's charter}
\label{sec:arc}

The spin sub-arc was planned as three papers, and the plan is recorded
verbatim in the arc's development notes:
\begin{itemize}[leftmargin=2em]
\item \textbf{Spin I}: given $\rout = 2\rin$, derive $L = \hbar/2$.
The captured-DP orbital geometry with the $1/r^2$ force law gives the
angular-frequency ratio $\omega_{\rm in}/\omega_{\rm out} = 2\sqrt{2}$
and total orbital angular momentum $\hbar/2$ exactly when
$\rin = \aBohr/[4(1+\sqrt{2})^2]$.
\item \textbf{Spin II}: derive $\rout = 2\rin$ from ZBW sub-harmonics.
The ZBW cloud is an open--closed radial resonator; its Mode~2 anchors
the inner $+$CP at the interior antinode ($\rth/3$) and the outer $-$CP
at the interior node ($2\rth/3$).
\item \textbf{Spin III}: derive the sub-harmonic spectrum from the
lattice Voronoi-cell eigenvalues --- ``the last mathematical gap to the
lattice.''
\end{itemize}
The honesty constraint on this third paper was also set in advance, in
the same notes: the lattice claim must not be asserted without either a
numerical eigenvalue computation of the lattice cell or an analytic
selection argument from the $2I$ symmetry group.  This V0 obeys that
constraint.  It consolidates what is provable now, diagnoses why the
existing numerical attempt did not settle the capstone, and freezes the
corrected instrument that can.

% =====================================================================
\section{The continuum spectrum, verified two ways}
\label{sec:spectrum}

\begin{proposition}[Resonator spectrum; Spin II]
The ZBW radial problem in the $\ell = 0$ sector, with a node at $r = 0$
(CP Exclusion) and a free antinode at $r = \rth$, has modes
\begin{equation}
\psi_n(r) = \sin(k_n r), \qquad
k_n = \frac{(2n-1)\pi}{2\rth}, \quad n = 1, 2, 3, \ldots
\label{eq:spectrum}
\end{equation}
\end{proposition}

\begin{theorem}[Mode-2 geometry; Spin II]
Mode 2 has its interior antinode at $\rin = \rth/3$ and its interior
node at $\rout = 2\rth/3$, whence $\rout/\rin = 2$ exactly ---
integer arithmetic of trigonometric zeros, independent of every
physical constant.
\end{theorem}

Table~\ref{tab:fd} reproduces the arc's committed finite-difference
verification (forward-difference Neumann ghost; the committed CSV was
regenerated bit-identically this session) alongside this paper's
independent reconstruction (symmetric lumped-mass Neumann formulation
--- a different discretization route).  Both agree with
Eq.~\eqref{eq:spectrum} at the $10^{-4}$ level across eight modes, and
the independent route locates the Mode-2 antinode and node at $R/3$ and
$2R/3$ to $0.1\%$ with ratio $2.000$ (verify script, checks 1--2).

\begin{table}[h]
\centering
\caption{Resonator eigenvalues: analytic vs.\ committed FD
(\texttt{data/spin3\_eigenvalues.csv}).  The independent lumped-mass
reconstruction agrees with both to $<0.05\%$
(\texttt{scripts/3234\_verify\_spin3\_v0.py}).}
\label{tab:fd}
\begin{tabular}{cccc}
\toprule
mode $n$ & $k_n$ (analytic) & $k_n$ (FD) & rel.\ error \\
\midrule
1 & 1.57079633 & 1.57063926 & $1.0\times10^{-4}$ \\
2 & 4.71238898 & 4.71191761 & $1.0\times10^{-4}$ \\
3 & 7.85398163 & 7.85319551 & $1.0\times10^{-4}$ \\
4 & 10.99557429 & 10.99447263 & $1.0\times10^{-4}$ \\
5 & 14.13716694 & 14.13574866 & $1.0\times10^{-4}$ \\
6 & 17.27875959 & 17.27702330 & $1.0\times10^{-4}$ \\
7 & 20.42035225 & 20.41829623 & $1.0\times10^{-4}$ \\
8 & 23.56194490 & 23.55956715 & $1.0\times10^{-4}$ \\
\bottomrule
\end{tabular}
\end{table}

% =====================================================================
\section{The scale-bridge identity}
\label{sec:bridge}

Spin~I places the inner orbit at the Bohr-derived scale
$\rin^{\rm (I)} = \aBohr/[4(1+\sqrt{2})^2]$; Spin~II places it at the
ZBW-cloud scale $\rin^{\rm (II)} = \rth/3$.  The committed script
records their ratio as $35.267491$ and observes that it equals
$6/[\alpha\cdot4(1+\sqrt{2})^2]$.  This paper upgrades that observation
to a theorem-grade statement:

\begin{theorem}[Scale bridge, exact]
Given the standard relations $\aBohr = \rC/\alpha$ and $\rth = \rC/2$,
\begin{equation}
\frac{\rin^{\rm (I)}}{\rin^{\rm (II)}}
= \frac{\aBohr/[4(1+\sqrt{2})^2]}{\rth/3}
= \frac{3\,\aBohr}{4(1+\sqrt{2})^2\,\rth}
= \frac{6}{4\alpha(1+\sqrt{2})^2}
\approx 35.2675
\label{eq:bridge}
\end{equation}
exactly --- an algebraic identity, not a numerical coincidence.
\end{theorem}

\noindent (Verified numerically to $10^{-12}$; verify check~3.)  The
physical reading, following the arc's development record: the
sub-harmonic constrains the \emph{ratio}, the Coulomb force balance
constrains the \emph{scale}, and the bridge between the two scales
involves $\alpha$ in exactly the combination \eqref{eq:bridge}.  The
mode structure and the orbital scale are two constraints, not one
constraint stated twice, and their compatibility is what
Eq.~\eqref{eq:bridge} records.

% =====================================================================
\section{The division of labor: selection, not corrections}
\label{sec:division}

The discreteness parameter at the ZBW scale is
\begin{equation}
\left(\frac{\lP}{\rth}\right)^2 \approx 7.0\times10^{-45},
\end{equation}
(verify check~3), so lattice discreteness cannot measurably shift the
continuum eigenvalues \eqref{eq:spectrum} at any accessible precision.
This sharpens what Spin~III must actually deliver: the lattice's job is
not to \emph{correct} the spectrum but to \emph{select} it --- to show
that the substrate's cell geometry, with the CP-Exclusion node and the
thermal free boundary as its physical boundary conditions, produces the
open--closed mode family (and in particular admits the Mode-2
node/antinode topology the captured DP anchors to) rather than some
other family.  A selection statement survives at any lattice fineness;
a correction statement would be empty at $10^{-45}$.

% =====================================================================
\section{The committed lattice-cell computation: what it did and did
not show}
\label{sec:instrument}

The arc's committed script (16 March 2026) contains, as its Part~3, a
lattice-cell computation: rejection-sample the interior of the 24-cell
(unit circumradius, the origin included as a sample point), build a
Gaussian-kernel $k$-nearest-neighbour graph, take the normalised graph
Laplacian $L_{\rm sym} = I - D^{-1/2}WD^{-1/2}$, and inspect the
radially averaged low eigenmodes for a ``Mode-2-like'' profile (exactly
one interior radial sign change).  Executed this session at both
reduced and default sampling ($8\,000$ and $20\,000$ points), the
computation reports, by its own diagnostic, \emph{no Mode-2-like mode}
--- and this V0's finding is that it cannot, for a stated reason:

\begin{proposition}[Boundary-condition mismatch]
The kNN graph Laplacian on uniformly sampled interior points
approximates the Laplacian with \emph{closed} (Neumann-type, free)
boundary behaviour everywhere; nothing in the construction encodes the
CP-Exclusion Dirichlet node at $r = 0$.  Its lowest nontrivial modes
are therefore the four 4D dipole modes of the closed cell, not the
open--closed radial family.
\end{proposition}

\noindent The diagnosis is verified reproducibly (verify check~4, same
construction, same seed): the lowest nontrivial spectrum is a 4-fold
near-degenerate group (spread $\sim10\%$, gap ratio $\sim2$ to the next
group --- the dipole signature in four dimensions), and no mode among
the lowest ten has exactly one interior radial zero.  Two consequences,
stated plainly:
\begin{itemize}[leftmargin=2em]
\item The Part-3 computation, as committed, neither confirms nor
disconfirms the capstone claim.  It solves a different boundary-value
problem.  (The committed figure, consistently, carries panels for
Parts~1, 2 and 4 only.)
\item Nothing upstream is affected: Parts~1--2 (the continuum spectrum
and its FD verification) are independent of Part~3 and reproduce
bit-identically.
\end{itemize}

\subsection{The domain question (founder ruling requested)}
\label{sec:domain}

The committed computation takes the \textbf{24-cell} as its domain.
The arc's charter language says ``600-cell Voronoi eigenvalue
spectrum.''  These are different objects, and V0 declines to paper over
the difference: whether the 24-cell is (a)~the intended Voronoi domain
of the CPP Grid-Point lattice, (b)~a deliberate tractable surrogate for
it, or (c)~a stand-in that should be replaced by the true GP-lattice
Voronoi cell, is a physical-picture question reserved to the founder.
It is registered as assumption A1 of the corrected instrument below and
flagged in \texttt{OPEN-QM-8}.

% =====================================================================
\section{The corrected instrument, frozen before its data exists}
\label{sec:corrected}

Per corpus practice (freeze the repair by principle, before the testing
data exists), the corrected lattice-cell measurement is specified now:

\begin{enumerate}[leftmargin=2em]
\item \textbf{Boundary conditions encoded, not implied.}  Dirichlet on
a small central core (all sample points with $r < \epsilon$ pinned to
zero --- the CP-Exclusion node), free (natural/Neumann) outer boundary
--- the open--closed problem the physics states.
\item \textbf{Radial-mode identification by isotropy score, not raw
sign counting.}  For each low eigenmode, compute the fraction of its
variance explained by its radial average; score modes with fraction
above a declared threshold as radial candidates; count interior zeros
on those only.  (The committed sign-counter applied to raw radial
averages of anisotropic modes is what produced zero counts of 13--36.)
\item \textbf{Domain per the founder's A1 ruling}: 24-cell, or the
GP-lattice Voronoi cell, or both as a robustness pair.
\item \textbf{Frozen readings.}  \textsc{mode2-recovered}: a radial
candidate with exactly one interior zero, node in $[0.60, 0.73]$ and
antinode in $[0.28, 0.40]$ of the boundary radius, stable across two
seeds and two sampling densities.  \textsc{mode2-absent}: no such mode
among the lowest twenty radial candidates at the highest density.
\textsc{indeterminate}: anything else $\Rightarrow$ instrument work,
not a verdict.
\item \textbf{Worker expectation, declared pre-run
(anti-extraction).}  \textsc{mode2-recovered} --- the continuum limit
is exact, discreteness corrections sit at $10^{-45}$, and the closed
problem's dipole modes are removed by the Dirichlet core; the risk is
angular interleaving in four dimensions, which the isotropy score is
designed to separate.
\end{enumerate}

This measurement, plus the analytic $2I$-symmetry selection argument
(the second route the arc's notes named), constitute
\texttt{OPEN-QM-8}.  Neither is attempted in V0.

% =====================================================================
\section{Epistemic ledger}
\label{sec:ledger}

\begin{itemize}[leftmargin=2em]
\item \textbf{Exact:} the resonator spectrum \eqref{eq:spectrum} and
Mode-2 ratio (Spin~II, restated and re-verified here two ways); the
scale-bridge identity \eqref{eq:bridge}; the discreteness bound.
\item \textbf{Established numerically:} FD eigenvalues to $10^{-4}$
across eight modes; Mode-2 geometry to $10^{-3}$; committed CSV
regenerated bit-identically.
\item \textbf{Diagnosed, not settled:} the lattice-cell selection
computation --- the committed instrument solved the closed problem;
capstone remains \texttt{OPEN-QM-8} with the corrected instrument
frozen above.
\item \textbf{Assumption awaiting founder ruling:} A1, the Voronoi
domain (Section~\ref{sec:domain}).
\item \textbf{Inherited:} everything Spin~I and Spin~II inherit; in
particular the $L = \hbar/2$ chain is complete at the continuum level
(Spin~I $+$ Spin~II) independently of this paper's open capstone.
\end{itemize}

% =====================================================================
\section*{Provenance}

Drafted at founder request (Session 149) from the arc's own materials:
the development-notes dialogue (Thomas Lee Abshier with Claude and
Grok, March 2026), the committed eigenvalue script and its data and
figure.  No prior Spin~III paper existed (full-history finding,
Patch 3232 record).  Verify:
\texttt{scripts/3234\_verify\_spin3\_v0.py} --- 10/10, including the
reproducible Part-3 diagnosis.

\begin{thebibliography}{9}
\bibitem{spin1} T.~L. Abshier, \emph{Emergent Spin-$\tfrac12$ from
Captured Dipole Particle Orbital Geometry} (spin sub-arc, Spin I).
\bibitem{spin2} T.~L. Abshier, \emph{Derivation of the Standing-Wave
Sub-Harmonic Condition for Captured Dipole Particle Orbits} (spin
sub-arc, Spin II).
\bibitem{zbwmass} T.~L. Abshier, \emph{ZBW $\hbar$ and Mass Units}
(SR companion c04) --- source of the thermal-boundary condition.
\end{thebibliography}

\end{document}
